\chapter{Storage, Indexing, and Persistence}
This chapter describes the data persistence and retrieval mechanisms that power the internship search platform. We detail the hybrid storage strategy that combines local file-based data for the internship catalog with cloud-based relational storage for user interactions, following established patterns in data integration and vector management \cite{doan2012principles, wang2021milvus}.

\section{Hybrid Data Storage Architecture}
The system adopts a hybrid approach to data management, balancing the high-speed read requirements of a large catalog with the transactional persistence needed for user-specific features \cite{doan2012principles}.

\subsection{Static Catalog and Data Normalization}
The primary repository for internship listings is a large-scale \texttt{intern\_data.json} file. To ensure data integrity, the system implements a robust ingestion pipeline:
\begin{itemize}
    \item \textbf{Data Cleaning:} Raw strings are normalized using techniques described in data cleaning literature \cite{chu2016data}.
    \item \textbf{Stable Identifier Generation:} The system generates deterministic IDs to enable effective record linkage and deduplication \cite{christen2012data, elmagarmid2007duplicate}.
    \item \textbf{Type Safety:} Ingested objects are transformed into a formal semantic schema \cite{sarawagi2008information}.
\end{itemize}

\subsection{Relational Persistence (Supabase)}
For dynamic user state, the system utilizes Supabase, a backend leveraging PostgreSQL. This relational foundation provides the necessary transactional support for features like user bookmarks and saved roles.

\section{Indexing and Metadata Enrichment}
Beyond simple storage, the system enriches the data to facilitate efficient discovery, particularly for high-dimensional semantic search \cite{johnson2019billion, wang2021milvus}.

\subsection{Heuristic Categorization}
As the source data is often unstructured, the system performs real-time metadata enrichment:
\begin{itemize}
    \item \textbf{Category Mapping:} Roles are classified into domains using keyword-based heuristics \cite{boselli2018classifying}.
    \item \textbf{Skill Extraction:} Descriptions are parsed against a predefined taxonomy to create a tag-based index \cite{giabelli2021skills, bhola2020retrieving}.
\end{itemize}

\subsection{Performance Optimization}
To handle the overhead of processing thousands of records, two layers of caching are implemented. This ensures that the system maintains sub-100ms response times even during heavy query loads \cite{buttcher2016information, menghani2023efficient}.

\section{Retrieval and API Design}
The system's data access layer is exposed via REST endpoints, designed for resilience and scalability.

\subsection{The Internship API}
The API serves as the unified data source, incorporating fail-safe logic to handle malformed input or missing catalog files.

\subsection{Search and Similarity Foundations}
The architecture includes utilities for similarity search. While current filtering relies on category matches, the system is structured to support future vector-based recommendation engines by leveraging Approximate Nearest Neighbor (ANN) indexing techniques such as HNSW \cite{malkov2018efficient, johnson2019billion}.
